\section{The approximate continuous model} \label{sec_wall_law}The purpose of this section is to analyze system \eqref{limit_equation}, and to prove Theorem \ref{thm2} and Corollary \ref{coro1}. We assume here that $f$ and $\rho$ are smooth and compactly supported in~$\overline{\Omega^0}$.

\subsection{Asymptotic expansion}
We will construct an approximation of the solution $(u^\eps,p^\eps)$ of \eqref{limit_equation}, of the following form:
\begin{equation*}
\left\{
\begin{aligned}
(u^\eps_{\app},p^\eps_{\app})(x_1,x_2,x_3) & \approx \sum_{i=0}^m \eps^i (u^i, p^i)(x_1,x_2,x_3), \quad x = (x_1,x_2,x_3) \in \Omega^0, \\
(u^\eps_{\app}, p^\eps_{\app}(x_1,x_2,x_3) & \approx \sum_{i=0}^m \eps^i (U^i, P^i)(x_1/\eps,x_2,x_3), \quad x = (x_1,x_2,x_3) \in \Omega^\eps\setminus\Omega^0.
\end{aligned}
\right.
\end{equation*}
Note that this approximation is made of two parts: an interior one, in $\Omega^0$, with a regular expansion in $\eps$, and a boundary layer part, localized in the depletion layer $\Omega^\eps \setminus\nobreak \Omega^0$. The boundary layer expansion involves boundary layer profiles $U^i = U^i(s,x_2,x_3)$, with the variable $s \in (-1,0)$ that stands for
$x_1/\eps$. It is also convenient to set:
\[
(u^i, p^i)= 0, \quad (U^i, P^i) = 0 \quad \text{ for } \: i < 0.
\]
Plugging the interior expansion in the Stokes equation, we~find: for all $i \ge 0$, in $\Omega^0$
\begin{equation} \label{stokesint}
\begin{aligned}
-\Delta u^i + \na p^i & = \delta_{0i} \big(f - \rho e \big), \\
\div u^i & = 0.
\end{aligned}
\end{equation}
Plugging the boundary layer expansion in the Stokes equation, we~find: for all $i \ge 0$, for all $(s,x_2,x_3) \in (-1,0) \times \R^2$:
\begin{align}
-\pa^2_s U^i_1 - (\pa^2_2 + \pa^2_3) U^{i-2}_1 + \pa_s P^{i-1} & = 0, \label{stokesx} \\
-\pa^2_s U^i_2 - (\pa^2_2 + \pa^2_3) U^{i-2}_2 + \pa_2 P^{i-2} & = 0, \label{stokesy}\\
-\pa^2_s U^i_3 - (\pa^2_2 + \pa^2_3) U^{i-2}_3 + \pa_3 P^{i-2} & = 0, \label{stokesz} \\
\pa_s U^i_1 + \pa_2 U^{i-1}_2 + \pa_3 U^{i-1}_3 & = 0 \label{stokesdiv}
\end{align}
The Dirichlet condition $u^\eps_{\app}\vert_{\pa \Omega^\eps} \approx 0$ further yields: for all $i \ge 0$,
\begin{equation} \label{dirichlet}
U^i\vert_{s = -1} = 0.
\end{equation}
Eventually, conditions
\[
[u^\eps_{\app}]\vert_{\pa \Omega^0} \approx 0, \quad [\Sigma(u^\eps_{\app},p^\eps_{\app})n]\vert_{\pa \Omega^0} \approx 0,
\]
which reflect continuity of the velocity field and the stress tensor at the interface $\pa \Omega^0$ give for all $i \ge 0$:
\begin{align}
U^i\vert_{s=0} & = u^i\vert_{x_1=0} \label{continuity_vel} \\
\pa_s U^i \vert_{s=0} - P^{i-1}\vert_{s=0} (1,0,0)^t & = \pa_1 u^{i-1}\vert_{x_1=0} - p^{i-1}\vert_{x_1=0} (1,0,0)^t. \label{continuity_stress}
\end{align}

\subsubsection*{Computation of the first terms.} We compute $U^0, u^0,p^0$. First, from \eqref{stokesdiv} and \eqref{dirichlet}, we~get $\pa_s U^0_1 = 0$ and $U^0_1\vert_{s=-1} = 0$, which implies $U^0_1 = 0$. Then, from \eqref{stokesy}, \eqref{dirichlet} and~\eqref{continuity_stress}, we~find
\[
\pa^2_s U^0_2 = 0, \quad U^0_2\vert_{s=-1} = 0, \quad \pa_s U^0_2\vert_{s=0} = 0,
\]
which leads to $U^0_2 = 0$. Similarly, $U^0_3 = 0$. Hence, $U^0 = 0$. Next, we~consider \eqref{stokesint} and~\eqref{continuity_vel} at rank $i=0$. They imply that $(u^0,p^0)$ satisfies \eqref{limit_equation}. In~particular, as~$f$ belongs to $L^{6/5}(\Omega_0) \cap H^m(\Omega^0)$, $(\na u^0,p^0)$ belongs to $H^m(\Omega^0)$ for all $m$.

\subsubsection*{Computation of next order terms.} We now assume that $i \ge 1$ and that profiles $U^k$, $P^{k-1}$, $u^k$, $p^k$ are known for all $k \le i-1$, with:
\[
U^k, P^{k-1} \in H^\infty((-1,0) \times \R^2), \quad \na u^k, p^k \in H^\infty(\Omega^0).
\]
We now show how to construct $U^i,P^{i-1}, u^i, p^i$. Considering \eqref{stokesdiv} and \eqref{dirichlet} yields
\[
U^i_1(s,\cdot) = - \int_{-1}^s (\pa_2 U^{i-1}_2 + \pa_3 U^{i-1}_3)(t,\cdot) dt \in H^\infty((-1,0) \times \R^2).
\]
Next, we~consider \eqref{stokesx} and \eqref{continuity_stress} which allow the calculation of \hbox{$P^{i-1} \!\!\in\! H^\infty(\!({-}1,0)\!\times\! \R^2)$}:
\[
P^{i-1}(s,\cdot) = \pa_s U^i_1\vert_{s=0} - \pa_1 u^{i-1}_1\vert_{x_1=0} + p^{i-1}\vert_{x_1=0}
+ \int_0^s \big(\pa^2_s U^i_1 + (\pa^2_2 + \pa^2_3) U^{i-2}_1 \big)(t,\cdot) dt.
\]
Next, we~consider \eqref{stokesy}--\eqref{dirichlet}--\eqref{continuity_stress}: these relations imply
\[
U^i_2(s,\cdot) = - \int_{-1}^s \int_{0}^t \Big((\pa^2_2 + \pa^2_3) U^{i-2}_2 + \pa_2 P^{i-2} \Big)(t',\cdot) dt' dt + (s+1) \pa_1 u^{i-1}_2\vert_{x_1=0}.
\]
A similar expression holds for $U^i_3$. Both $U^i_2, U^i_3 \in H^\infty((-1,0) \times \R^2)$. Eventually, we~see by \eqref{stokesint} and \eqref{continuity_vel} that $(u^i,p^i)$ solve a homogeneous Stokes equation with the inhomogeneous Dirichlet condition
\[
u^i\vert_{x_1=0} = U^i\vert_{x_1=0}.
\]
By standard regularity results for the Stokes equation, \cf
\cite[Th.\,IV.3.2 \& IV.3.3]{Galdi}, $\na u^i, p^i \in H^\infty(\Omega^0)$.

